\documentclass[twoside]{article}
\usepackage[utf8]{inputenc}

% The package supports other formats as well, but this template
% is primarily checked and updated for examinations at TUM BGDM
%\usepackage[uebung, answers]{tumbgdm}

\usepackage[uebung]{tumbgdm}
\usepackage{pythonhighlight}

\Professor{Prof. Dr. Martin Werner}
\Universitaet{Technical University of Munich}
\Institut{School of Engineering and Design}
\Lehrstuhl{Professorship Big Geospatial Data Management}
\Semester[WiTe 2024/25]{Winter Term 2024/25}
\Vorlesung{Principles of Programming}
\InstitutsWebsite{https://www.bgd.lrg.tum.de/}
\VorlesungsWebsite{https://www.moodle.tum.de/course/view.php?id=103704}

\nummerUebungsblatt{3}
\nummerErsteFrage{1}

\LearningOutcome{ Python Packages: Numpy, Pyplot}

\begin{document}
\maketitle

\begin{task}{Numpy}{}{}
%\begin{lstlisting}[numbers=none]
%import numpy as np
%\end{lstlisting}
Import the NumPy package (Official documentation: \href{https://numpy.org/doc/}{https://numpy.org/doc/}) and use it to solve the following tasks:

\begin{enumerate}
    \item Print the Numpy version and configuration
    \item Create a 10x10 matrix filled with zeros
    \item Find the shape of an arbitrary array (e.g. the one from a.)
    \item Create a 4x6 matrix with zeros but the value at index [2, 3] which is 1
    \item Create a 3x7 matrix filled with values ranging from 21 to 42
    \item Create a vector of size 10 filled with random values from 0 to 1
    \item Reverse the vector from f. (first element becomes last)
    \item Create a 3x3 matrix with random integer values ranging from 0 to 8
    \item Find indices of non-zero elements from [1,0,7,0,1,0,0,1,0,0,1]
    \item Find all unique values from [1,0,7,0,1,0,0,1,0,0,1]
    \item Create a 3x3 identity matrix
    \item Multiply a 5x3 matrix by a 3x2 matrix
    \item Create a random vector of size 10 and sort it
    \item Create a 5x5x5 array with random values and find the minimum, the maximum and the mean value
\end{enumerate}
\end{task}

\clearpage
\begin{task}{Matplotlib}{}{}
%\begin{lstlisting}[numbers=none]
%from matplotlib import pyplot as plt
%\end{lstlisting}
Import the Pyplot package from Matplotlib (Official documentation: \href{https://matplotlib.org/stable/api/index.html}{https://matplotlib.org/stable/api/index.html}) and use it to solve the following tasks:

\begin{enumerate}
    \item Load an image of your choice
    \item Display the image
    \item What are the spatial dimensions, data type, and range of the values of the image?
    \item What is the color value of the pixel at position [21,42]?
    \item Calculate the mean colour value for each channel over all pixels
    \item Convert the image in grayscale and display it. The gray value h can be calculated from the RGB values using the formular h = 0.2989r + 0.5870g + 0.1140b.
    \item Change the colormap so that the image is actually shown as grayscale and display it again
    \item Save the grayscale image in the PNG-format
    \item Create and plot a histogram of the grayscale image
\end{enumerate}
\end{task}

\begin{task}{Numpy continued}{}{}
Use Numpy to solve the following tasks.
\begin{enumerate}
    \item Create a 2D array with 1 on the border and 0 inside
    \item Normalize a 5x5 random matrix
    \item Multiply a 4x3 matrix by a 3x5 matrix
    \item Create a random vector of size 6 and replace the maximum value by 0
    \item Create a 2D integer matrix and return the transpose
    \item Create a checkerboard 10x10 matrix using the "tile" function
\end{enumerate}
\end{task}


\clearpage
\begin{task}{Matplotlib continued}{}{}
Use Matplotlib to solve the following tasks.
First, we create some random dataframe with pandas:
\begin{lstlisting}[numbers=none]
import pandas
df = pandas.DataFrame(data={"a":np.random.randint(0, 100, 50),
                            "b":np.random.randint(0, 100, 50),
                            "c":np.random.randint(0, 100, 50)})
\end{lstlisting}

\section{Line Plot}
\begin{enumerate}
    \item Create a figure of size 15x8 with two subplots, top and bottom
    \item Draw two lines (data a and b) in the top plot, one blue and one orange, and add a grid
    \item Create two vectors of size 50 filled with random values in range 0-10 and use "matplotlib.pyplot.fill\_between" function to plot an artificial standard deviation for both lines (Hint: Set "alpha=0.2" and "edgecolor=None")
    \item For the top plot, add a title and a legend in the top-middle
    \item In the bottom plot, create a graph of only the data points (c), marked by circles, but with no line connecting the points. Further, add a grid, title, and legend for the bottom plot
\end{enumerate}

\section{Bar Plot}
\begin{enumerate}
    \item Create a figure of size (15,18) with three subplots, aligned vertically
    \item Create bar graphs for data a and data b in the first two subplots but in third subplot put data b on top of data a
    \item Add a legend and title for each plot
    \item Plot the magnitude of each bar on top of it
\end{enumerate}

\end{task}

\end{document}